\section{Introduction}
Tinnitus and hearing loss are among the most common auditory health concerns
worldwide, affecting communication, sleep, mood, work participation, and quality
of life. More than 1.5 billion people live with some hearing loss, a number
projected to rise by 2050 \citep{who2021world}, and tinnitus affects an estimated
10--15\% of adults, a meaningful subset of whom experience bothersome tinnitus
\citep{tunkel2014cpg,mccormack2016review}. In employed middle-aged adults,
tinnitus frequently emerges alongside psychosocial stress, sleep disruption,
noise, and delayed help-seeking; because these factors are intertwined, simple
causal narratives are unreliable and can mislead patients and clinicians.

\subsection{Motivating observation and naming principle}
Consider an anonymized composite trajectory: a worker in their late 40s develops
unilateral tinnitus during heavy workload, which resolves over months; about a
year later the other ear develops sudden fullness and reduced audibility, but---
unaware of the narrow treatment window for sudden sensorineural hearing loss
(SSNHL)---the person delays care, and recovery is incomplete. Two clinically
actionable points follow: stress-associated tinnitus is common and often
self-limiting, but sudden or rapidly progressive hearing loss is an urgency in
which delayed care, not stress, drives irreversible impairment. Attributing acute
hearing symptoms to ``stress'' without assessment is hazardous. Accordingly, STARS
(Stress, Tinnitus, and Audiometric Research Study) adopts a strict naming
principle: because KNHANES and NHANES carry general perceived-stress and mood
items rather than validated job-strain instruments, we never call the exposure
measured ``occupational stress.'' It is labeled \emph{perceived stress}, with
\emph{work-related factors} (occupation, employment, and where available hours and
shift work) modeled as distinct covariates.

\subsection{A survey-weighted study reported against a prespecified plan}
This is a survey-weighted study whose \emph{primary} hypothesis test---the
perceived-stress analysis in KNHANES---is reported in full
(Section~\ref{sec:results}), with supporting NHANES analyses. To limit analyst
degrees of freedom, the design, variable definitions, cycles, adjustment sets, and
evaluation plan were fixed \emph{before} analysis, and committed code and outputs
make every value reproducible. Three further components---cross-national
validation, clinician-verified AI extraction, and a clinical extension---are
\emph{prospective} (design only) and kept subordinate to the empirical finding.

\subsection{Aims and hypotheses}
Aim~1 is the empirical contribution reported here; Aim~3 is developed \emph{with
results} in the companion SAFE-EAR paper, Aim~4's open pipeline accompanies this
article, and Aim~2 is prospective.
\begin{enumerate}[leftmargin=1.5em,itemsep=2pt]
\item \textbf{Association (development).} Estimate weighted associations between
perceived stress and (a) tinnitus and (b) audiometric hearing loss in KNHANES,
adjusting for noise and health covariates. \emph{H1: a positive stress--tinnitus
association persists after adjustment, while the association with audiometric
thresholds is weaker and may attenuate to non-significance after adjusting for
noise and age.}
\item \textbf{Transportability.} Test whether harmonized models developed in
KNHANES retain discrimination and calibration in NHANES. \emph{H2: discrimination
degrades modestly but calibration is restorable by recalibration.}
\item \textbf{Explainable, safe AI (companion paper).} Evaluate whether
schema-constrained extraction feeding a deterministic red-flag layer routes SSNHL
and neurologic warning signs to urgent referral. \emph{H3: near-complete red-flag
recall, bounded by extraction quality.} Developed with results in SAFE-EAR.
\item \textbf{Open, shared benefit.} Provide an openly licensed, reproducible
pipeline (code, data dictionary, model cards, prompts).
\end{enumerate}

\subsection{Contributions}
Primary: a survey-weighted, public-data result showing that perceived stress
dissociates the tinnitus \emph{symptom} from the audiometric \emph{threshold},
reported without unsupported causal claims and reproducible end-to-end. Secondary:
an open, reproducible scaffold (harmonized mappings, design-based estimators,
model cards, safety tests) that lets others regenerate every number and add local
cohorts. The remaining components---a prespecified KNHANES$\rightarrow$NHANES
external-validation plan and a clinically governed AI/referral-safety component in
which a deterministic red-flag layer overrides probabilistic predictions---are
prospective (design only); any prediction model is a research risk-stratification
tool, never a screening, diagnostic, or triage device.
